% ============================================================
%  Section 4: Dataset
%  H&E to Sirius Red Virtual Staining Paper
% ============================================================

\section{Dataset}
\label{sec:dataset}

The dataset introduced in this work is, to our knowledge, the first publicly
released resource dedicated to unsupervised H\&E~$\to$~Sirius Red translation
research in liver tissue.
The H\&E and SR images originate from the same liver specimens described in
\textcolor{red}{[TODO: cite ASBTi paper --- Ghallab et al., JHEP Reports 2025,
doi:10.1016/j.jhepr.2025.101599]}, where they were used for histopathological
quantification of hepatic fibrosis.

% -------------------------------------------------------
\subsection{Animal Model and Tissue Collection}
\label{sec:animal_model}
% -------------------------------------------------------

Obstructive cholestasis was induced in eight-to-ten-week-old male C57BL/6N mice
(Janvier Labs, France) by ligation of the extrahepatic common bile duct~(BDL) at
a position between the gallbladder and the duodenum, following previously
established procedures~\textcolor{red}{[TODO: cite original BDL protocol paper]}.
Sham-operated control mice underwent the identical surgical procedure without
bile-duct ligation.
Both BDL and sham animals received vehicle as part of a larger study evaluating
pharmacological ASBT inhibition~\cite{ghallab2025asbt}; only the vehicle-treated
arms are included in this dataset.
The model recapitulates the progressive hepatic fibrosis observed in cholestatic
liver disease~\cite{bataller2005liver}, with fibrosis severity increasing systematically
across four post-surgical time points: day~3 (early stage, minimal fibrosis),
day~21 (intermediate fibrosis with inflammation), day~42 (established fibrosis),
and day~63 (advanced fibrosis).
Animals were housed under standard environmental conditions with free access to
water and ad libitum feeding on a standard rodent diet (Ssniff, Soest, Germany).
All experiments were ethically approved by the local committee~(LANUV,
North Rhine-Westphalia, Germany; application number 81-02.04.2022.A286).

In total, 36 animals contributed liver sections to this dataset, spanning four
post-surgical time points and both treatment conditions~(BDL+vehicle and
sham+vehicle), with four to five animals per cohort.

% -------------------------------------------------------
\subsection{Histological Staining}
\label{sec:staining}
% -------------------------------------------------------

Consecutive 4\,$\mu$m-thick sections were cut from paraformaldehyde-fixed,
paraffin-embedded~(PFPE) liver tissue blocks, ensuring that adjacent sections
sample the same tissue volume with minimal inter-section deformation.
The first section of each pair was stained with haematoxylin and eosin~(H\&E)
using the Discovery Ultra Automated Slide Preparation System
(Ventana Medical Systems / Roche).
The immediately adjacent section was stained with Sirius Red using a
commercially available kit according to the manufacturer's
instructions\textcolor{red}{[TODO: specify kit name/manufacturer if available]}.
Sirius Red binds selectively to fibrillar collagens (types I and III) through
electrostatic interactions between the sulphonic acid groups of the dye and the
basic amino acids of collagen, providing high-specificity collagen
visualisation~\cite{constantine1969}.
Because both sections are cut from the same PFPE block, macroscopic tissue
morphology is consistent across the H\&E and SR stains; however, microscopic
deformation between consecutive sections introduces a degree of spatial
misalignment that precludes exact pixel-level correspondence
(addressed in Section~\ref{sec:metrics}).

% -------------------------------------------------------
\subsection{WSI Acquisition}
\label{sec:wsi_acquisition}
% -------------------------------------------------------

All stained slides were digitised on a Zeiss Axio Scan.Z1 whole-slide scanner at
$\times$20 magnification, corresponding to a native pixel resolution of
$0.221\,\mu$m/pixel.
Images were acquired in CZI format and exported to RGB TIFF using
QuPath~\cite{bankhead2017qupath} prior to tiling.
Slides with severe tissue artefacts (folds, air bubbles, or scanning failures
affecting the majority of the section) were excluded prior to tiling.
Slides exhibiting a mild yellow background tinge in non-collagen tissue areas ---
a minor staining variability that does not affect PSR-positive signal ---
were retained in the training set, as this appearance is present in the real SR
distribution and is accounted for by the downstream segmentation model.

% -------------------------------------------------------
\subsection{Tiling and Preprocessing}
\label{sec:dataset_tiling}
% -------------------------------------------------------

Tissue masks were generated for each WSI using a pixel classifier in
QuPath~\cite{bankhead2017qupath}: tissue regions were identified by applying an
intensity threshold on the Eosin channel for H\&E slides and on the DAB channel
for SR slides.
The masked WSI was tiled into non-overlapping $256{\times}256$\,pixel patches
using a stride equal to the tile size.
Tiles in which less than 50\% of pixels fell within the tissue mask were discarded
from the training set; no tissue filtering was applied to the test set.
All retained tiles were normalised channel-wise to $[-1,\,1]$ by subtracting
$0.5$ and dividing by $0.5$ per channel.
This preprocessing was applied identically to both the H\&E~(domain~A) and
SR~(domain~B) tile sets.

Table~\ref{tab:tile_counts} reports the resulting tile counts per domain and split.

\begin{table}[h]
\centering
\caption{Tile counts after tissue filtering. Training subsets (25\%, 50\%) are
strict subsets of the full 100\% training set.}
\label{tab:tile_counts}
\begin{tabular}{lcccc}
\toprule
\textbf{Split} & \textbf{WSIs} & \textbf{H\&E tiles} & \textbf{SR tiles} & \textbf{Data fraction} \\
\midrule
Train (100\%) & 30 & 192{,}847 & 183{,}253 & 100\% \\
Train (50\%)  & 15 & 109{,}049 &  96{,}292 & 50\%  \\
Train (25\%)  &  7 &  49{,}312 &  49{,}293 & 25\%  \\
Test          &  5 &  16{,}285 &  16{,}285 & ---   \\
\bottomrule
\end{tabular}
\end{table}

% -------------------------------------------------------
\subsection{Data Splits and Subsets}
\label{sec:splits}
% -------------------------------------------------------

The dataset was partitioned at the \textit{animal level} to prevent data leakage:
no WSI from the same animal appears in more than one split.
Splits were stratified by time point and treatment group to ensure that all
disease stages and both conditions~(BDL, sham) are represented in every
training subset and in the test set.
The training set comprises 30 WSIs organised into sequentially numbered per-WSI
subdirectories (\texttt{001/} through \texttt{030/}) for each domain, enabling
reproducible data-fraction subsets via contiguous folder ranges:
25\%~$(\llbracket 1,7 \rrbracket)$, 50\%~$(\llbracket 1,15 \rrbracket)$, and
100\%~$(\llbracket 1,30 \rrbracket)$.
Smaller subsets are strict subsets of larger ones, ensuring that performance
differences across data fractions are attributable solely to training set size and
not to variation in which slides are included.
The test set comprises five held-out WSIs from animals not present in any training
fraction, and is shared identically across all 54 experimental configurations.

Table~\ref{tab:split_stats} summarises the per-split slide counts and disease-stage
distribution.

\begin{table}[h]
\centering
\caption{Dataset split summary. BDL stage indicates post-ligation day at which
tissue was harvested; sham = control mice without bile-duct ligation.}
\label{tab:split_stats}
\begin{tabular}{lcccccc}
\toprule
\textbf{Split} & \textbf{Animals} & \textbf{WSIs} &
\multicolumn{4}{c}{\textbf{BDL stage (WSIs)}} \\
\cmidrule(lr){4-7}
& & & Sham & Day 3 & Day 21/42 & Day 63 \\
\midrule
Train & 28 & 30 & \textcolor{red}{[TODO]} & \textcolor{red}{[TODO]} & \textcolor{red}{[TODO]} & \textcolor{red}{[TODO]} \\
Test  &  5 &  5 & \textcolor{red}{[TODO]} & \textcolor{red}{[TODO]} & \textcolor{red}{[TODO]} & \textcolor{red}{[TODO]} \\
\bottomrule
\end{tabular}
\end{table}

% -------------------------------------------------------
\subsection{Public Release}
\label{sec:public_release}
% -------------------------------------------------------

The full dataset --- including all H\&E and SR WSIs, tissue masks, and tile-level
metadata --- will be released publicly under a
\textcolor{red}{[TODO: licence, e.g.\ Creative Commons Attribution 4.0
(CC~BY~4.0)]} licence alongside the publication of this paper.
The dataset will be deposited at \textcolor{red}{[TODO: repository, e.g.\ Zenodo /
BioImage Archive]} and assigned a persistent digital object identifier~(DOI)
\textcolor{red}{[TODO: DOI once assigned]}.
Pre-trained model checkpoints and all training scripts will be released in the same
repository to enable full reproduction of the reported results.


% ============================================================
%  New bibliography entries for this section
%  (merge with references.bib)
% ============================================================
%
% @article{ghallab2025asbt,
%   author  = {Ghallab, Ahmed and Myllys, Maiju and Gonz{\'a}lez, Daniela and
%              others},
%   title   = {Stage-dependent effects of systemic {ASBT} inhibition in a
%              cholestasis-induced cholemic nephropathy mouse model},
%   journal = {JHEP Reports},
%   volume  = {7},
%   year    = {2025},
%   doi     = {10.1016/j.jhepr.2025.101599}
% }
%
% @article{constantine1969,
%   author  = {Constantine, V. S. and Mowry, R. W.},
%   title   = {Selective staining of human dermal collagen.
%              {II}. {T}he use of {Picrosirius} Red {F3BA} with polarization microscopy},
%   journal = {Journal of Investigative Dermatology},
%   volume  = {50},
%   number  = {6},
%   pages   = {419--423},
%   year    = {1969},
%   doi     = {10.1038/jid.1968.68}
% }
%
% (Shared entries already listed:
%  bataller2005liver, huang2013sirius)
